\section{Problem formulation and security definitions}
\label{sec:problem}

\subsection{State and transitions}

The trusted store maps each identifier $n$ to a record
\begin{equation}
\Omega[n]=(s,H,a,z,r),
\label{eq:record}
\end{equation}
where $s$ is the value, $H=(n,q,\sigma,d,\tau,t_i,t_e,u,c)$ is metadata, $a=F_K(\mathsf{enc}(H))$, $z\in\{\mathsf{live},\mathsf{reserved},\mathsf{consumed},\mathsf{expired},\mathsf{revoked}\}$ is lifecycle state, and $r$ is an optional reservation nonce. The visible handle is $h=(n,q,a)$ plus an explicitly non-sensitive label. Storing the complete $H$ server-side prevents the agent from editing a destination even if the display representation is malleable.

The abstract machine supports:
\begin{align}
  (h,\Omega') &\leftarrow \Mint(s,q,\sigma,d,\tau,c,t_i,t_e,u;\Omega), \label{eq:mint}\\
  (y,\Omega') &\leftarrow \Resolve(h,x,\tau',t,\gamma;\Omega), \label{eq:resolve}\\
  O' &\leftarrow \Redact(O,R), \label{eq:redact}
\end{align}
where $x$ is the TCB-derived target context, $\gamma$ is confirmation evidence, $y$ is either an internal injection capability or $\Reject$, $O$ is a raw observation, and $R$ is the active redaction registry. Expiry and revocation are state transitions that never return $s$.

\begin{definition}[Target match]
For trusted context $x=(p,o,f,v)$, $\Match(d,x)=1$ exactly when every mandatory component of destination predicate $d$ accepts the corresponding component in $x$. An unavailable mandatory identity produces $0$, not a wildcard.
\end{definition}

\begin{definition}[Eligible handle]
A record is eligible at time $t$ for task $\tau'$ and target $x$ if its authenticator verifies, its state is live, $t_i\le t<t_e$, $u>0$, $\tau'=\tau$, $\Match(d,x)=1$, and confirmation evidence satisfies $c$.
\end{definition}

\subsection{Observation boundary}

The agent view $\View_A(\pi)$ contains agent-visible screenshots, accessibility trees, handle metadata, action results, timing rounded to the public granularity, and application messages. It excludes TCB memory, sealed storage, protected-confirmation pixels, and the private input buffer between resolver and target. This boundary is intentionally concrete: claiming that ``the agent never sees the value'' without enumerating observations would be unfalsifiable.

Let $\ell(s,H,\pi)$ be allowed leakage. By default it contains $q$, a coarse destination label, issue age bucket, remaining-lifetime bucket, success/failure status, and effects already observable from the task, such as arrival at a logged-in screen. It must not include $s$, a deterministic hash of a low-entropy OTP, value length unless unavoidable, or error codes that distinguish candidate values.

\begin{definition}[Agent-view non-disclosure]
For equal-length secrets $s_0,s_1$ with the same allowed leakage and any probabilistic polynomial-time adversary $\mathcal{A}$ controlling the agent, the distinguishing advantage is
\begin{equation}
\begin{aligned}
\Adv^{\mathsf{nd}}_{\SH,\mathcal{A}}(\lambda)=
\big|&\Pr[\mathcal{A}(\View_A(\Real(s_0)))=1]\\
&-\Pr[\mathcal{A}(\View_A(\Real(s_1)))=1]\big|.
\end{aligned}
\label{eq:nd}
\end{equation}
Non-disclosure holds when this advantage is negligible given equal $\ell$.
\end{definition}

This is a non-interference-style definition with declassified effects. It is stronger than ``the prompt omits the value'' but weaker than preventing inference from all world state. For example, a successful login reveals that the hidden code was accepted; that outcome is intentionally declassified.

\begin{definition}[Destination confinement]
For any trace, if \Resolve\ returns an injection capability containing or referencing $s$ for target $x$ under task $\tau'$, then $\Match(d,x)=1$ and $\tau'=\tau$ for the live record bound to $h$.
\end{definition}

\begin{definition}[Single-use replay safety]
For a handle minted with $u=1$, at most one resolve operation reaches the committed-injection state, even when an adversary schedules an arbitrary number of concurrent calls.
\end{definition}

\subsection{Correctness and availability}

Security is vacuous if the resolver always rejects. Functional correctness requires that an eligible handle aimed at a supported, stable target is injected exactly once unless a platform error occurs, and that the failure is surfaced without revealing $s$. Availability is not absolute: the OS, user, or app may cancel; the target may disappear; a confirmation UI may be unavailable. The evaluation therefore separates policy rejection, platform failure, agent failure, and user denial.

\subsection{Research questions}

The evaluation is organized around six questions:

\begin{description}[leftmargin=1.5cm,style=nextline]
\item[RQ1] Does OS-source minting keep raw values out of all instrumented agent observations and traces?
\item[RQ2] Does destination/task binding prevent injection-driven redirection without blocking legitimate completion?
\item[RQ3] Do lifecycle and atomic reservation prevent replay, expiry bypass, and concurrency races?
\item[RQ4] Does observation redaction prevent post-entry rereading in pixels and accessibility trees?
\item[RQ5] What latency, memory, energy, and user-interaction overheads arise by source and handle type?
\item[RQ6] Which identity ambiguities, unsupported apps, WebViews, keyboards, or OEM modifications cause failure or force a weaker policy?
\end{description}
